\chapter{认知查询与检索}

\section{认知查询}

认知查询表达调用方的语义意图，而不暴露具体检索算法开关。查询可以指定主体、领域和对象范围，以及实体、时间、语义线索、硬约束、软偏好、调用方能力、计算预算、材料展开要求和诊断需求。

硬约束决定对象是否合法；软偏好只影响合法对象之间的排序。按稳定身份进行的精确读取直接访问权威对象，不经过相似度检索。

\section{查询计划}

查询执行前生成固定查询计划（QueryPlan）。计划在看到候选结果之前确定检索通道、各通道预算、融合参数、拓扑预算、可选重排、材料展开预算和一致性要求。

计划中的某个检索通道没有产生结果时，该通道贡献为零。结果集合本身不得反向改变其他检索通道的权重或归一化尺度。

\section{检索通道}

\begin{ReportWideTable}{>{\raggedright\arraybackslash}p{3.0cm}>{\raggedright\arraybackslash}p{10.4cm}}
\captionof{table}{认知检索通道}\label{tab:retrieval-lanes}
\toprule
\textbf{通道} & \textbf{职责}\\
\midrule
精确读取 & 按对象身份、修订和精确键读取。\\
实体检索 & 按认知指向、实体绑定和来源提及召回相关对象。\\
词项检索 & 处理短语、代码、路径、专名和其他词项线索。\\
稠密检索 & 按规范表示、语义片段或来源约束的召回视图检索。\\
时间检索 & 按发生、观察、有效、形成和记录时间过滤或排序。\\
运行时检索 & 从活动认知、驻留集和工作上下文中取得当前相关对象。\\
拓扑检索 & 在有界关联拓扑中生成候选排序。\\
资源检索 & 查询外部知识或数据源，并按需将实际接触结果形成观察事件。\\
可选重排 & 对已通过硬约束的候选集执行二次排序。\\
\bottomrule
\end{ReportWideTable}

\section{候选身份与多视图}

候选的语义身份由主体、领域所有者、对象身份和修订身份确定。一个修订可以拥有规范表示、段落表示、召回线索和其他服务视图；这些视图属于同一认知对象，不构成多份独立证据。

每个检索通道先将视图命中聚合到对象修订层，再参与跨通道融合。视图数量不得直接形成额外投票权。

\section{融合与最终校验}

跨通道融合函数由查询计划固定，并满足以下约束：参数在候选生成前确定；缺失通道的贡献明确；候选集合中出现无关对象时，不改变既有通道的定义；拓扑检索以普通通道身份参与，不使用计划之外的隐藏加分。

融合与可选重排完成后，返回结果前重新核验当前修订、对象生命周期、访问策略、访问控制、来源状态以及抑制和清除状态。服务投影过期时可以降级、重建或明确失败，不能把过期投影当作权威认知。

\section{诊断}

查询结果能够区分权威对象不存在、访问受限、对象被抑制、服务投影不可用或过期、预算不足、来源不可展开、线索不足和真正未知。预算导致的召回不足不得被报告为主体没有相应认知。
